% Generated draft: requires mathematical and source review.
\documentclass[11pt]{article}
\usepackage[margin=25mm]{geometry}
\usepackage{fontspec}
\setmainfont{Latin Modern Roman}
\usepackage{amsmath,amssymb,mathtools,unicode-math}
\setmathfont{Latin Modern Math}
\usepackage{enumitem,needspace,xurl,hyperref}
\setlength{\parindent}{0pt}
\setlength{\parskip}{5pt}
\setlength{\emergencystretch}{3em}
\newcommand{\N}{\mathbb{N}}
\newcommand{\Z}{\mathbb{Z}}
\newcommand{\Q}{\mathbb{Q}}
\newcommand{\R}{\mathbb{R}}
\newcommand{\C}{\mathbb{C}}
\newcommand{\F}{\mathbb{F}}
\newcommand{\A}{\mathbb{A}}
\newcommand{\PP}{\mathbb{P}}
\newcommand{\E}{\mathbb{E}}
\newcommand{\eps}{\varepsilon}
\newcommand{\dd}{\,\mathrm{d}}
\begin{document}
\section*{UnsolvedMath ID: 20003131}
\textbf{Source problem code:} catalogue entry
\subsection{Definitions and mathematical statement}
Let \(X\) be a finite-dimensional positively curved Alexandrov space and \(Y\subset X\) a totally geodesic subspace of codimension \(k=\dim X-\dim Y\ge1\). In the usual compact, closed-subspace interpretation, the source asks whether
\[
H_i(X\setminus Y;\Z)=0\qquad\text{for all }i>2k-2.
\]
The source does not explicitly specify compactness, boundary, coefficients, connectedness or the precise intrinsic/extrinsic meaning of totally geodesic in the singular setting. Integral homology is an explicit formulation; other coefficients or reduced homology should be stated separately. In particular codimension one would allow nontrivial \(H_0\) while requiring positive-degree vanishing.

The supplied aid separates two inputs behind the smooth result. For a closed oriented \(n\)-manifold \(X\) and a suitable closed locally contractible subspace \(Y\), Alexander--Lefschetz duality gives
\[
H_i(X\setminus Y;R)\cong H^{n-i}(X,Y;R).
\]
If the relative groups on the right vanish in degrees \(j\le n-2k+1\), the desired cutoff follows. Wilking's \((n-2k+1)\)-connected inclusion supplies such relative vanishing in its manifold setting. For nonorientable manifolds the integral statement uses the orientation local system; \(\Z/2\) avoids that issue. In singular Alexandrov spaces neither the inclusion-connectivity assertion nor this manifold duality can simply be assumed. A valid extension must supply the relevant relative-cohomology vanishing and an appropriate ambient duality, or another argument.

The aid proves a singular join special case. For a spherical join \(X=A*B\) with \(Y=A\) the indicated factor, the complement strongly deformation-retracts onto \(B\), by increasing the join parameter to the \(B\)-end. Since
\[
k=\dim B+1,
\qquad H_i(X\setminus A;R)\cong H_i(B;R)=0\quad(i>k-1),
\]
this gives a stronger cutoff whenever the factors satisfy the Alexandrov/join hypotheses and homological dimension properties used in the construction. Those full factor hypotheses are not reproduced here.

The aid emphasizes the example \(X=\operatorname{Susp}(\C\PP^m)=\C\PP^m*S^0\), \(Y=\C\PP^m\), with the appropriately scaled positively curved metric, \(m\ge1\). This equator can obstruct a naive extension of Wilking's inclusion-connectivity conclusion, but its complement consists of two contractible components. It therefore satisfies the codimension-one complement bound and is not a counterexample to the source question. The general Alexandrov complement bound remains open in the supplied account.
\subsection{Short English statement}
Ask whether removing a codimension-\(k\) totally geodesic subspace leaves homology only through degree \(2k-2\). The supplied reduction separates connectivity from duality, proves a stronger join-factor case, and explains why a projective-suspension connectivity counterexample does not refute the complement bound.
\subsection{Sources}
\begin{itemize}
\item Problem 6, AIM workshop \emph{Manifolds with nonnegative sectional curvature}, \url{https://aimath.org/WWN/nnsectcurvature/nnsectcurvature.pdf}. The source names the Wilking Connectivity Theorem but supplies no full bibliography.
\item The research summary invokes Alexander--Lefschetz duality and the projective-suspension test, without separate full bibliographic entries or complete singular-space hypotheses.
\item Local UnsolvedMath registry, permanent ID 20003131, exact code \texttt{catalogue entry}; canonical location \texttt{aim-topology-notes.json}, notes[218], dated 14 August 2026; statement labeled exact.
\item Research attempt 1: partial result, problem open at the run, rigor label proved, novelty confidence low; registry status partially solved. The duality reduction and join special case are machine-generated research aids, not peer reviewed; independent verification is required before citation.
\end{itemize}
\end{document}
